% Editable transcription; verify reused text against the original PDF.
\section{Is AI intelligent?}

Although these recent breakthroughs are very impressive and based on
methods with parallels in the human brain, the question whether such
systems display real intelligence remains relevant.

On this, opinions are divided.
As said, AI already outperforms humans in several areas.
In other areas, more computation power, a combination of existing
techniques and new techniques that are under development now, may
lead to similar levels of (super)human performance.
According to Newell and Simon's (2007) physical symbol system
hypothesis, there is no principal reason why an extrapolation of
recent advances in AI should not lead to true intelligence.

Yet, the general assessment of how intelligent AI really is, and how
much it still differs from human intelligence, is contradictory,
at best.
On the one hand, cognitive psychologists and neuroscientists generally
acknowledge the merits of modern AI methods (Naselaris et al., 2018),
especially in the context of perception (VanRullen, 2017) and language
(Henderson, 2020).
On the other hand, others argue that previous critiques of early AI
systems---including the famous arguments by Searle (1980),
Dreyfus (1965), and Penrose (1994)---still hold for modern AI systems.
Common to these critiques is that AI systems lack understanding.

Chess again helps to illustrate the problem.
In certain so-called ``anticomputer positions'' (such as the
impenetrable fortress of pieces shown in Fig. 2), human chess players
quickly see the point of the position, while AI systems are only able
to evaluate the position after considering long sequences of moves.
For example, the Leela Chess Zero engine has a difficult time with
this kind of position.
Humans somehow have the ability to use both general reasoning and
automatic pattern recognition when playing chess, which suggests
that the critique about lack of understanding is still relevant.

In the case of chess, the ``lack of understanding'' limitation is
perhaps temporary.
Leela chess might, in time, learn to understand ``anti-computer
positions'', where perhaps a modular approach in which Leela chess
is combined with a symbolic or Bayesian approach might work.
But this clearly would not end the debate.
According to Searle, Dreyfus and others, the problem lies deeper.
AI systems lack semantics, feelings, a body and goals.
Some of these limitations can perhaps be overcome by using robots
that do have a body or real time interactions with the actual
(rather than digital, simulated) world, which is often thought
to be necessary to grasp causality---an element that is sorely
absent from most AI systems
(Lake et al., 2017; Tenenbaum et al., 2011).

\begin{quote}
\small
\textbf{Figure 2 (image omitted).}
A typical anti-computer position.
Human chess players quickly see that black, in spite of its material
advantage, can not make progress.
The best computer chess programs assess the position as much
better for black.
\end{quote}

So how should we assess the ``lack of understanding'' issue from
the intelligence researcher's point of view?
We will answer this with two definitions of intelligence in mind.
The first is a practical definition: intelligence is what the
intelligence test measures.
This definition is not taken very seriously
(however, see van der Maas, Kan, \& Borsboom, 2014), but how well AI
performs on human intelligence tests provides an interesting insight
into the limitations of AI nonetheless
(which is discussed in the next section).

In the most commonly accepted definitions of intelligence, we see
two important requirements.
One is the ability to deal with various forms of information and
to solve all kinds of cognitive problems, and the other is the ability
to both quickly and effectively learn how to deal with new situations
(Legg \& Hutter, 2007).

AI clearly meets the first requirement, as there is arguably nothing
wrong with AI's crystallized intelligence: AI programs can play games
like go and chess, compose music, create art, solve complex math
problems, provide medical diagnosis, etc.
But things are different when we confront a trained AI system with
completely novel situations.
Schank (1980) already mentioned the problem of generalizability.
Humans are able to transfer solutions from one problem domain
to another, apply general, abstract concepts (such as presented
in ``anticomputer problems'' in chess) in reasoning and to develop
solutions for completely new problems.
AI, on the other hand, often fails to learn such abstract rules and
if they do, it requires an immense amount of training examples
(Lake et al., 2017).

We think it is justified to say that generalization is still a
weakness of AI systems.
At the same time, we see many promising developments.
We already mentioned transfer learning, which works by virtue of
learning general representations that facilitate broad generalization.
Multitask learning is an active area of research (Ruder, 2017).
Ongoing attempts to integrate deep learning and symbolic AI or
Bayesian reasoning may also be crucial for progress in this
research area.
New work on one-shot learning
(Vinyals, Blundell, Lillicrap, Kavukcuoglu, \& Wierstra, 2016)
may also lead to breakthroughs.

Whether these future developments also make AI systems `understand'
things like we do is a topic of hot philosophical debate
(for further discussion we refer to Cole, 2020).

\subsection{How well do AI systems perform on intelligence tests?}

Newell (1973) argued that if a single AI system could solve a diverse
set of intelligence test problems, then we could consider it to be
intelligent.
Since the 1960's there have been waves of efforts to have AI systems
solve intelligence test problems.
This started with Evan's (1964) ANALOGY program that could solve
geometric analogies from the WAIS and was intended to be able
to generalize to other reasoning tasks (but not pursued).
In Hernández-Orallo, Martínez-Plumed, Schmid, Siebers, and Dowe's
(2016) review we see that AI has tackled a broad range of intelligence
subtests, ranging from letter and number series to block design and
exclusion tasks, and without a doubt excels at what are considered
tests of crystallized intelligence (e.g., math, vocabulary) and
processing capacity (e.g., memory, speed).
The one test that has continually popped up in the literature
over the years is Raven's Progressive Matrices (RPM).
AI programs have been used to study how people solve the RPM
(Carpenter, Just, \& Shell, 1990; Lovett \& Forbus, 2012) to fully
independent programs that take raw images as input and induce rules
and select the correct option
(Correa, Prade, \& Richard, 2012; Kunda, McGregor, \& Goel, 2013)
or even generate it (Pekar, Benny, \& Wolf, 2020), in many cases
reaching human level performance.

An especially interesting case is that of Bongard problems
(Bongard, 1967 in Hofstadter, 1979).
In Fig. 3 you see a simple Bongard problem where the test-taker has
to induce which pattern the images on the left follow and those
on the right do not.
Hofstadter (1979) considered these types of reasoning tasks to lie
at the core of human intelligence.
His PhD student Foundalis (2006) created Phaeaco, a
(by today's standards) rudimentary AI system, to try to solve
Bongard problems.
Phaeco contains a cognitive architecture with modules to process
the images, match visual patterns, learn, and to store and retrieve
learned patterns.
Phaeco was relatively successful, although the program did not reach
human level performance.
More recently, Nie et al. (2020) tried to solve Bongard problems
using state-of-the-art deep learning methods.
Their best model achieved 66\% accuracy, and yet performance still
falls behind that of humans who are able to solve about 90\% of
their Bongard benchmark problems.

So, although in many cases intelligence tests can be solved by AI
programs, most of these are specialized programs that can only solve
one particular task.
Some attempts have been made to create more general AI systems.
For example, in Nie et al.'s (2020) study one of the models they used
to solve Bongard problems was built upon one of the most successful
Raven-like item solvers
(Barrett, Hill, Santoro, Morcos, \& Lillicrap, 2018), alas this model
performed only around the level of chance.
So currently, we cannot say that AI is ``intelligent'' by the
psychometric definition in the sense of being able to achieve human
level performance on all subtests of an intelligence test using
one general AI system.

But, what if we approach this question from the other side:
How well do humans perform on tests of intelligence for AI?
Over the AI lifespan, an interesting field has emerged that takes
the psychometric approach to assessing general intelligence in AI,
but then applied to problems generated specifically for AI research
(e.g., Bringsjord \& Schimanski, 2003; Chollet, 2019;
Evans, 1964; Hernández-Orallo et al., 2016).
This approach is characterized by hundreds or thousands of
(often computer generated) items fit for a non-verbal test of
intelligence (Barrett et al., 2018; Chollet, 2019;
Liu et al., 2019; Nie et al., 2020;
Zhang, Gao, Jia, Zhu, \& Zhu, 2019).
These item banks are referred to as benchmarks and used to objectively
compare the performance of different AI systems, and, in some cases,
also compare AI performance to that of humans.
There is a clear role for Intelligence researchers in this field
given our expertise in test development, sampling from the general
population and in assessing human performance.
There is also a role for AI in improving existing measures of human
intelligence (e.g., automatic item generation;
Gierl \& Haladyna, 2012).
For AI, currently more intelligence test benchmarks are needed
(e.g., for problem analogies, Ichien, Lu, \& Holyoak, 2020).
Also, these benchmarks only sometimes include humans among their
comparison samples, which appear limited to adults
(Nie et al., 2020) or only a few experts (Chollet, 2019).
Unsurprisingly, at this point in time humans generally outperform
AI systems on these types of AI ``fluid'' intelligence tests
(e.g., Chollet, 2019; Nie et al., 2020).
Given the activity in the field, perhaps many strides will be made
in the coming years to achieve human-like performance on these
benchmarks, but also psychology can inspire innovation in AI by
providing a developmental account of if and how humans learn
to solve these AI benchmarks.

\begin{quote}
\small
\textbf{Figure 3 (image omitted).}
Bongard visual pattern recognition problem, where the test-taker
has to induce which pattern the images on the left follow and those
on the right do not.
In this example, the images on the left are convex and those on
the right are concave.
From
\url{https://upload.wikimedia.org/wiki/commons/thumb/e/e9/BongardProblem_convex_poligons.svg}.
\end{quote}
